\hypertarget{appendix-2a2f1b74bdd3}{}
\section*{Terminating fractions and dyadic rationals}
\section{Terminating fractions and dyadic rationals}\label{terminating-fractions-and-dyadic-rationals}

\textbf{Curriculum depth:} Introduced · \textbf{First seen:} C005 · \textbf{Latest development:} C005

\subsection{Definition and mechanism}\label{definition-and-mechanism}

A reduced rational p/q terminates in base B iff q divides a power of B. In binary q must be a power of two. Binary32 0.1 is exactly 13421773/2\^{}27, chosen as an approximation of non-dyadic 1/10.

\subsection{Prerequisites}\label{prerequisites}

\begin{itemize}
\tightlist
\item
  \hyperref[appendix-83c3f28f7bda]{Positional notation} --- Finite fractional digits correspond to an integer over a power of the base.
\item
  \hyperref[appendix-5057f9ad327f]{Binary32} --- Its finite significand/exponent makes each stored finite value a dyadic rational.
\end{itemize}

\subsection{Used by}\label{used-by}

Connect later dependent concepts here when their lessons arrive.

\subsection{Recall and next depth}\label{recall-and-next-depth}

Reconstruct the definition's example and name the assumptions that make it valid.

\textbf{Pending:} Other bases, decimal conversion algorithms and exact rational arithmetic.

\subsection{Source and navigation}\label{source-and-navigation}

\href{Continuous\%20Notes/Computing\%20Foundations\%20-\%20From\%20Bits\%20to\%20C++\%20and\%20Systems\%20APIs\#When\%20arithmetic\%20meets\%20uncertainty}{M005 teaching chapter} · \hyperref[appendix-fd090d982a15]{Part definitions} · \href{}{Dependency Map (02-Foundations-II, appendix-4ad47bdd269a)}

\subsection{Why decimal-friendly numbers can be binary-hostile}\label{why-decimal-friendly-numbers-can-be-binary-hostile}

A reduced rational number \(p/q\) has a terminating representation in base \(b\) exactly when every prime factor of \(q\) also divides \(b\). Base 10 has prime factors 2 and 5, so fractions such as \(1/10\) terminate. Base 2 has only prime factor 2, so only \textbf{dyadic rationals} with denominator \(2^k\) terminate.

That is why \(0.5\), \(0.25\), and \(0.125\) are exact in ordinary binary floating point while \(0.1\) repeats forever and must be rounded. The issue is not that floating point is ``bad at decimals''; it is a consequence of positional representation.

This perspective helps debugging: multiplying a binary approximation to 0.1 by ten need not recreate the exact path one imagines from decimal arithmetic. For currency or protocol fields whose semantics are decimal-exact, scaled integers or decimal floating formats may be better contracts.

A reconstruction exercise should reduce several fractions, classify their denominator prime factors, and predict whether they terminate in bases 2, 10, and 16 before checking actual floating encodings.
